\section{Introducción y motivación}

El sector espacial atraviesa un cambio de paradigma que se conoce como \textbf{NewSpace}: frente
al modelo clásico de misiones únicas, de plazos de una década y presupuestos de cientos de
millones, se impone el de constelaciones de satélites pequeños, fabricados en serie, con ciclos
de desarrollo de meses y costes dos órdenes de magnitud menores. Ese cambio no es solo
económico, sino técnico, y afecta directamente a cómo se diseña la electrónica de a bordo. La
vía tradicional ---componentes cualificados para espacio, endurecidos frente a radiación y con
herencia de vuelo demostrada--- es incompatible con esos plazos y esos costes: el catálogo de
piezas cualificadas es reducido, caro y va tecnológicamente por detrás. El enfoque NewSpace
acepta por ello el uso de \textbf{componentes comerciales} (COTS, \textit{commercial off-the-shelf}),
trasladando la fiabilidad desde la pieza individual hacia la arquitectura del sistema: redundancia,
detección de errores y capacidad de recuperación en vuelo. Esta filosofía es la que enmarca todas
las decisiones de diseño de este trabajo, y explica que la plataforma de partida sea un MPSoC
comercial y que los componentes de las tarjetas diseñadas sean de grado industrial o automotriz
y no de grado espacial.

Dentro de ese marco, el diseño de ordenadores de a bordo (OBC) para satélites exige resolver un
reto de ingeniería que va más allá del procesamiento puro: conseguir que el procesador central se
comunique de forma determinista y fiable con una familia heterogénea de periféricos ---sensores de
actitud, actuadores, cargas de pago--- a través de buses serie con especificaciones eléctricas y de
temporización muy diferentes entre sí. En el marco del proyecto LINCE, el grupo B105
Electronic Systems Lab de la Universidad Politécnica de Madrid asume la responsabilidad del
subsistema de comunicaciones del OBC.

El trabajo se desarrolló en el laboratorio B105 en calidad de colaborador en prácticas dentro
del proyecto LINCE, y la tarea inicial consistió en implementar el soporte de comunicaciones
serie RS422 y RS485 sobre la plataforma de evaluación Zynq UltraScale+ ZCU102. Esta
plataforma, basada en un MPSoC que integra en un mismo encapsulado un sistema de
procesamiento ARM (PS) y una lógica programable tipo FPGA (PL), proporciona la
flexibilidad necesaria para implementar transceptores serie a medida. Se optó por un diseño
propio en VHDL en lugar de emplear IPs de terceros, con el objetivo de obtener control total
sobre el comportamiento del transceptor: parámetros de temporización, gestión de errores y
adaptación a los requisitos concretos del sistema.

El alcance del trabajo, sin embargo, no se mantuvo en ese encargo inicial, y la forma en que se
amplió merece ser explicada porque condiciona la estructura de esta memoria. Una vez operativo
el transceptor en la lógica programable, validarlo exigía ejercitar sus catorce canales sobre
buses reales, y la plataforma de evaluación no ofrece ese hardware. Se abordó entonces, por
iniciativa propia y sin que mediara encargo, el \textbf{diseño de una tarjeta de expansión propia}
que reprodujera las topologías RS-485 multipunto y RS-422 punto a punto sobre las que el sistema
tendría que operar en vuelo. El resultado de esa tarjeta fue lo que llevó al consorcio a asignar
al autor el diseño de las \textbf{dos tarjetas oficiales del proyecto} ---la placa CDHS, de gestión
de datos y control térmico, y la placa AOCS, de control de actitud y órbita---, con las que Sener
debía validar su propio firmware sobre la ZCU102. El trabajo pasó así de ser exclusivamente de
codiseño hardware-software a incluir también el diseño electrónico, la selección de componentes y
la fabricación de tres tarjetas de circuito impreso.

Esa ampliación fue posible porque la colaboración con Sener ---socio tecnológico del proyecto
LINCE--- fue continua y no puntual: el trabajo se integró en su ciclo de \textit{sprints}, con
reuniones periódicas de revisión y planificación que permitían reasignar tareas conforme
avanzaba el proyecto. A su vez, Sener traslada los requisitos técnicos impuestos por Indra, que
lidera el consorcio.

El resultado de todo este trabajo es una plataforma funcional que abarca el codiseño
hardware-software del subsistema de comunicaciones serie, su integración bajo el sistema
operativo de tiempo real RTEMS, y la fabricación y validación de las tres tarjetas de
prueba: las dos oficiales del proyecto (CDHS y AOCS) y la de comunicación serie de diseño
propio, caracterizada eléctricamente sobre sus catorce transceptores.

\section{Metodología}

El desarrollo de este trabajo se enmarcó en la dinámica de trabajo colaborativo del proyecto
LINCE, que sigue una \textbf{metodología ágil}. El trabajo se organizó en \textit{sprints}, con
reuniones quincenales de revisión y planificación con los ingenieros de Sener asignados al
proyecto ---donde se demostraba el incremento alcanzado, se identificaban bloqueos y se
comprometía el alcance del siguiente ciclo---, además de comunicación continua por correo para
resolver dudas técnicas y alinear requisitos. Las tareas se gestionaron mediante Jira como
tablero común. Esta cadencia es la que explica que el alcance del trabajo pudiera reajustarse
sobre la marcha, incorporando el diseño de las tarjetas CDHS y AOCS a mitad de desarrollo. Los
requisitos de diseño de las PCBs de prueba llegaron a través de Sener, que los traslada desde
Indra como líder del consorcio.

El trabajo se estructuró en tres fases principales:

\textbf{Fase 1 --- Familiarización y establecimiento del entorno (octubre -- enero):} Los primeros
meses se dedicaron a establecer las bases del entorno de desarrollo: instalación y
configuración de Vivado y Vitis, compilación de RTEMS 7 desde código fuente para la
arquitectura AArch64 mediante el RTEMS Source Builder, y validación del flujo completo
(síntesis en Vivado $\rightarrow$ empaquetado en Vitis $\rightarrow$ ejecución en RTEMS sobre la ZCU102) con
un ejemplo funcional de extremo a extremo. Esta fase fue fundamental para comprender la
arquitectura PS-PL del MPSoC y las particularidades del entorno de compilación cruzada
para sistemas empotrados.

\textbf{Fase 2 --- Desarrollo e integración (febrero -- junio):} Con el entorno establecido, se abordó
el desarrollo principal: el transceptor serie en VHDL, el driver para RTEMS, el diseño de las
PCBs CDHS y AOCS, y las aplicaciones de prueba. Se siguió una metodología iterativa: cada
módulo se implementaba, se validaba de forma aislada (mediante simulación en Vivado o
prueba directa sobre la placa) y se integraba en el sistema global antes de avanzar al
siguiente.

\textbf{Fase 3 --- Estudio del transporte y caracterización (junio -- julio):} La validación con los
catorce canales activos reveló que el factor limitante del sistema no era el transceptor sino el
transporte de datos entre el PS y la PL, lo que reorientó la fase final del trabajo: se
implementaron y compararon las tres arquitecturas de transporte (AXI GPIO, catorce AXI DMA y AXI
MCDMA con puente VHDL) sobre un banco de pruebas común, y se fabricó y caracterizó eléctricamente
la placa de comunicación serie de diseño propio mediante una campaña de nueve tests sobre sus
catorce transceptores.

Dos principios atraviesan todo el desarrollo y conviene enunciarlos aquí, porque explican
decisiones que aparecerán repetidamente a lo largo de la memoria.

El primero es la \textbf{verificación como parte del diseño, no como fase posterior}. Ningún bloque
de la lógica programable se integró en el sistema sin haber pasado antes su propio
\textit{testbench} en simulación, y el sistema completo se validó en placa mediante baterías de
pruebas automatizadas que actuaban como test de regresión. Este principio no es retórico: el
código de \textit{testbench} escrito en el trabajo (4.551 líneas) es prácticamente equivalente en
volumen al código VHDL sintetizable que verifica (4.710 líneas). Fue además determinante en el
punto más difícil del proyecto ---la depuración del bus AXI contra el motor de DMA---, donde la
ausencia de observabilidad desde el software hizo de la simulación la única vía de diagnóstico.

El segundo es la \textbf{automatización y la reproducibilidad del flujo}. Un ciclo de iteración en
esta plataforma atraviesa síntesis en Vivado, empaquetado en Vitis, compilación cruzada de RTEMS
y generación de la imagen de arranque; hecho a mano es lento y, sobre todo, propenso a errores
silenciosos. Se invirtió por ello un esfuerzo deliberado en construir herramientas propias
---\textit{scripts} TCL que regeneran el diseño de bloques completo desde cero, \textit{scripts} de
compilación y despliegue de la imagen de arranque, utilidades de instrumentación--- de modo que
cualquier resultado de esta memoria pueda reproducirse desde el repositorio con un único comando.
El peso de esa decisión es medible: las herramientas de automatización suponen el 40\,\% de las
líneas de código escritas en el trabajo, más que ninguna otra categoría.
